As in STAR and POPSTAR, we simulate client inputs by drawing from a Zipfian distribution with support $\zipfsup = 10{,}000$ and exponent $\zipfexp = 1.03$.
For our main experiments we set $\pd = 350$.
Then, given a report count $\max$ and threshold $\threshold$, we use the parameter selection algorithm described in \cref{subsubsec-distribution-optimizations} to find the values of $\np_i$ and $\batchsize_i$ for each step.
We use a field size equal to the smallest prime greater than $\max$ for each experiment.

We select $\pd = 350$ in order to make the closest comparison between our protocol and POPSTAR for a client count of $\max = 100,000$ and threshold $\threshold = 1000$.
Recall from \cref{subsec:popstar-leakage} that the most serious leakage in POPSTAR comes from the \emph{exact counts} that it leaks.
Under the parameters benchmarked in \cite{\citepopstar}, POPSTAR is expected to leak the exact count for $19$ values.
We therefore set $\pd = 350$, as it is the lowest value of $\pd$ for which there are expected to be $19$ values reported between $\pd$ and $\threshold = 1000$ times.



\subsection{Report Generation}

All report generation benchmarks were run on an AWS {\sf c7i.8xlarge} EC2 instance.


\begin{figure}[h]
	\centering
	\begin{tabular}{ |c|c|c|c|c| }
		\hline
		\makecell{Modulus}         &
		\makecell{Client\\Runtime} &
		\makecell{Server\\Runtime} &
		\makecell{Client\\Message} &
		\makecell{Server\\Message}                                              \\
		\hline
		$2048$                     & $1.50$s & $0.23$s & $35.47$ kb & $0.42$ kb \\
		$3072$                     & $2.83$s & $0.74$s & $53.13$ kb & $0.42$ kb \\
		\hline
	\end{tabular}
	\caption[Report generation runtimes]{Report generation runtimes for $\np = 139$ and $\pd = 349$. \emph{Modulus} indicates the length in bits of the modulus used for the Benaloh encryption scheme. \emph{Client Message} is the size of the message sent from each client to the Randomness Server. \emph{Server Message} is the size of the Randomness Server's response. All runtimes are single-threaded.}
	\label{fig:protocol-runtime}
\end{figure}

The benchmarks for our report generation protocol can be seen in \cref{fig:protocol-runtime}.
The benchmarks presented are for simultaneously evaluating $c = 138$ polynomials with degree $350$, which represents the largest parameters among our experiments in \cref{subsec:agg-experiments}.

Our protocol is efficient, with client and server computation on the order of seconds, and bandwidth in the range of $50$ kilobytes.
The benchmarks shown in \cref{fig:protocol-runtime} are all single threaded.
When up to $4$ threads are used, client computation drops to $0.39$ and $0.74$ seconds for the $2048$ and $3072$-bit moduli, respectively.
The bulk of the work done by the client is the homomorphic evaluation of the polynomials.
For the $2048$-bit modulus a single evaluation takes $6.4$ms, and so the full $138$ evaluations represent $59\%$ of the client's total runtime.

Our protocol is noticeably lighter over-the-wire than POPSTAR, which requires $363$kb of communication between the client and the Randomness Server.
However, our protocol has a higher computational cost, in POPSTAR clients are estimated run in $2$ milliseconds, though we view this tradeoff as acceptable given our stronger security guarantees.


\subsection{Aggregation}
\label{subsec:agg-experiments}

All benchmarks in this section were performed on a server with an AMD EPYC 7R13 64-Core Processor, 252 Gigabytes of RAM, and an RTX 3090 GPU.

We benchmark the aggregation algorithm on input sizes $10^5 \leq \max \leq 10^6$, set the threshold $\threshold$ to be $1\%$ of $\max$, and keep $\pd$ fixed at $350$.
The results of these experiments can be see in \cref{fig:gpu-main-aggregation-runtimes}.

\begin{figure}[h]
	\centering
	\begin{tabular}{ |c|c|c| }
		\hline
		$\max$          & \makecell{Aggregation\\ Runtime} & \makecell{Leaked Values\\(Expected)} \\
		\hline
		$100{,}000$     & $1$ hour, $10$ minutes           & $19$                                 \\
		\hline
		$250{,}000$     & $1$ hour, $12$ minutes           & $63$                                 \\
		\hline
		$500{,}000$     & $54$ minutes                     & $133$                                \\
		\hline
		$750{,}000$     & $1$ hour, $17$ minutes           & $202$                                \\
		\hline
		$1{,}000{,}000$ & $57$ minutes                     & $239$                                \\
		\hline
	\end{tabular}
	\caption{Aggregation runtimes and leaked values across client counts with $\threshold = \frac{\max}{100}$.}
	\label{fig:gpu-main-aggregation-runtimes}
\end{figure}




Due to the optimizations from \cref{sec:star-optimizations}, the runtime of our aggregation algorithm remains relatively constant at approximately $1$ hour.
The downside of this is that, as we keep $\pd$ fixed, the expected number of leaked values climbs as the number of total clients increases.

As compared to STAR and POPSTAR, our aggregation algorithm takes significantly longer: for $100{,}000$ clients STAR and POPSTAR aggregate in $2$ and $10$ seconds, respectively.
We note that this disparity reduces for larger client pools.
At 1 million clients STAR aggregates in $5$ minutes while POPSTAR takes $20$.

As it pertains to leakage, for all values of $\max$ greater than $100{,}000$ the number of values leaked by our scheme is greater than the number of exact counts leaked by POPSTAR (which stays fixed at $19$).
We believe that this could be an acceptable tradeoff for applications where guaranteeing protection for rare values, especially against a spying Aggregation Server, is paramount.


\subsection{Challenges in Increasing $\pd$}

All benchmarks in this section were performed on a DGX Spark \cite{nvidia_dgx_spark}, a consumer-focused device with a combined CPU+GPU system-on-a-chip and $128$GB of unified memory.

Optimally, $\pd$ would be nearly equal to $\threshold$, granting privacy for as many non-heavy-hitters as possible.
We investigate the feasibility of such configurations in this section by showing how the resources required by the aggregation algorithm, including the necessary VRAM, scale as $\pd$ approaches $\threshold$ for an instance with $\max = {50},000$ and $\threshold=1000$.
This can be seen in \cref{fig:increasing-tpriv}.

\begin{figure}[h]
	\centering
	\begin{tabular}{ |c|c|c|c| }
		\hline
		$\pd$ & \makecell{Aggregation\\ Runtime} & VRAM      & \makecell{Leaked Values\\(Expected)} \\
		\hline
		$350$ & $7$ minutes                      & $0.6$GB   & $10$                                 \\
		\hline
		$500$ & $14$ minutes                     & $0.8$GB   & $5$                                  \\
		\hline
		$600$ & $22$ minutes                     & $1.2$GB   & $4$                                  \\
		\hline
		$700$ & $43$ minutes                     & $2.1$GB   & $2$                                  \\
		\hline
		$800$ & $1$ hour, $34$ minutes           & $4.7$GB   & $1$                                  \\
		\hline
		$900$ & $2$ hours, $19$ minutes          & $18.82$GB & $1$                                  \\
		\hline
	\end{tabular}
	\caption{Resource usage and leaked values for an instance with $\max = 50{,}000$ and $\threshold = 1000$.}
	\label{fig:increasing-tpriv}
\end{figure}


We see in \cref{fig:increasing-tpriv} how increasing $\pd$ leads to sharp increases in resource use, approximately doubling the aggregation runtime every $100$ it increases.
The VRAM usage is even more dramatic.
Indeed, attempting to go one step further and aggregate over an instance with $\pd = 950$ would require $77$GB, well beyond the capacity of most consumer GPUs,  to construct the initial matrix.
It is this rapid scaling that prevents feasibly closing the gap between $\pd$ and $\threshold$.

